\section{System Architecture}\label{sec:system}
\subsection{Construction, Execution, and Reporting}
Figure~\ref{fig:architecture} separates responsibilities that can otherwise
be conflated in a model-generated report. A browser or command-line interface
collects the question. An external agent may submit the same structured world.
A ready world, user specification, or optional author
produces a proposal. Review fixes the participants, rules, information links,
selection mode, seeds, and execution budget. The engine then obtains decisions
and applies Algorithm~\ref{alg:round}. These are logical responsibilities within
one application, not a requirement for separately deployed services.

\begin{figure*}[t]
\centering
\begin{tikzpicture}[node distance=7mm and 8mm]
\node[box] (user) {Research question\\Browser, CLI, or agent client};
\node[box,right=of user] (author) {Ready world or author\\Participants and assumptions};
\node[box,right=of author] (review) {Review and validation\\Rules, scope, and budget};
\node[box,below=of review] (engine) {Synchronous engine\\Observe, choose, settle};
\node[box,left=of engine] (jev) {Typed decision adapter\\Jev or offline utility rule};
\node[box,left=of jev] (evidence) {Timestamped evidence\\Sources and known gaps};
\node[box,below=of engine] (record) {Recorded trajectories\\Replay and presentation};
\node[box,left=of record] (near) {Optional NEAR registry\\Artifact commitment};
\draw[arr] (user)--(author);\draw[arr] (author)--(review);
\draw[arr] (review)--(engine);\draw[arr] (engine)--(jev);
\draw[arr] (jev.north) to[bend left=15] (engine.north);
\draw[arr,dashed] (evidence)--(author);\draw[arr] (engine)--(record);
\draw[arr,dashed] (record)--(near);
\end{tikzpicture}
\caption{Boyun's execution architecture. Dashed links are optional. Evidence
informs the reviewed world; it does not automatically determine behavioral
coefficients. The registry binds a recorded artifact, rather than certifying
the simulated outcome.}
\label{fig:architecture}
\end{figure*}

Approval is associated with the reviewed configuration so that later edits
require a new review. Requests, responses, completed trajectories, and failures
are retained separately. Reports replay the recorded decisions before computing
metric paths, means, medians, and spending summaries. Grouped scenarios summarize
completed paths rather than adding unexecuted model-generated futures. The
report includes requested and completed run counts; otherwise failures could
silently change the denominator behind a chart.

\begin{table}[t]
\caption{Records retained for interpretation and reproduction.}
\label{tab:artifacts}
\centering\footnotesize
\begin{tabularx}{\columnwidth}{@{}p{0.28\columnwidth}Y@{}}
\toprule Record & Why it is needed \\
\midrule
Evidence snapshot & Fix the source, observation time, network, and coverage available to the analysis. \\
Reviewed world & Preserve assumed participants, resources, and consequences independently of inference. \\
Observations and decisions & Preserve what each participant was shown and selected; no internal reasoning trace. \\
Provider receipt & Record returned model identity, reported token use, and request/response fingerprints. \\
Trajectory and failures & Recompute settlement and retain incomplete outcomes in reporting. \\
Presentation & Explain recorded paths and trade-offs without requiring raw-log inspection. \\
Commitment opening & Allow a later consistency check against an on-chain registration. \\
\bottomrule
\end{tabularx}
\end{table}

\subsection{Evidence and Persona Grounding}
The public-data collector accumulates pool observations and available recent
trades, preserving failures, source cooldowns, and coverage gaps. NEAR and
Solana observations remain separate because their trading environments need
not imply the same behavior. Forward collection does not reconstruct an
unobserved 30-day wallet history.

Grounding proceeds from observable features to reviewed hypotheses: preserve
a timestamped snapshot, compute features such as trade sizes and frequency,
record unknowns, then propose participant roles for review. A large trade
supports a statement about observed size; it does not establish that the
wallet is a market maker, an insider, or one of several wallets controlled by
the same person. Likewise, news and social posts require source attribution,
timestamps, and deduplication. The current archive and wallet analysis do not
constitute a validated estimator of population preferences or causal effects.

\subsection{Distinct Execution Models}
The strategic engine uses cash costs, bounded public metrics, joint effects,
and subscribed publications. A separate financial laboratory represents
integer cash and token balances, fees, and a synthetic constant-product AMM.
In such an AMM, trades change the pool reserves and therefore the executable
price; the marked value of a holding can differ from the cash obtained by
selling it. Analytical components for finite games, Bayesian updating,
forecast scoring, and conditional event chains are also available separately.

These distinctions matter for interpretation. Not every strategic run invokes
every analytical component. In particular, a creator-profit question requires
sales, inventory, costs, liabilities, and withdrawals. The strategic engine's
utility and attention metrics do not provide that accounting. Combining the
strategic and financial models requires explicit economic transitions, rather
than relabeling a metric as dollars.

\subsection{Typed Inference and Resource Controls}
The Jev adapter uses TypeSafe AI's direct System One interface
\cite{typesafeapi}. It transmits the participant's state and typed questions,
then validates the returned answers and token usage. A versioned model identity
is recorded because an alias may change over time. Invalid responses cannot
settle a round. Authentication, billing, rate-limit, and service failures stop
execution; automatic inference retries and silent provider fallback are disabled.
A timed-out request can still have consumed remote resources, so an uncertain
outcome is retained rather than assumed free.

Normally there is one request per participant per round per world, except for
forced wait actions. This counting rule is needed to review the cost of a run
before dispatch. TypeSafe lists Jev 1.14 input at \$0.042 per million tokens
with free output tokens \cite{typesafemodels}. For illustration, 100 participants
over two rounds in one world require at most 200 actor requests. At an assumed
2,500 input tokens per request, the input-only estimate is \$0.021. This is a
planning example; observation length, authoring, hosting, and failed requests
change the complete operating cost. A durable call allowance limits dispatch
count, rather than guaranteeing a dollar ceiling.

\subsection{Agent Interoperability}\label{sec:mcp}
The Model Context Protocol (MCP) is a client--server interface through which
an AI application discovers tools, reads resources, and obtains reusable
prompts \cite{mcp2025}. It standardizes access to a system's capabilities;
it does not define the language model's reasoning policy or validate a
scientific result. An external agent using MCP is also distinct from a
simulated participant inside a Boyun world.

Boyun exposes a local standard-input/output server over the existing
construction and execution interfaces. Resources describe the world contract,
guidelines, and examples. A construction prompt guides the connected agent's
chosen model in supplying a world. Tools inspect the game construction and each own-information view, validate
that world, execute a bounded
offline run, inspect its status, and retrieve replay-verified results and a
Markdown report. The agent supplies data rather than code or arbitrary file
paths. Opaque run identifiers and bounded artifact reads limit the interface
to its explicit workspace.

The server's participant policy is the local immediate-utility rule, not Jev
inference. It exposes no credential import, paid inference, wallet signing,
or chain-submission tool. Its purpose is to let research agents construct
and inspect reproducible experiments through a common interface. A local
transport creates no network listener, but retrieved world data can still
reach a remote model through the connected agent application. MCP support
therefore does not establish end-to-end confidentiality. Compatibility is
evaluated with the versioned protocol specified in the integration experiment.

\subsection{Privacy and Artifact Registration}\label{sec:near}
Bring-your-own-key operation lets a user supply their own provider credential.
Session-only storage avoids deliberately persisting that credential in the
application, but both the host and selected provider can process plaintext
requests. A self-hosted installation still sends requests remotely when it
uses a remote model. Offline deterministic operation uses a local rule based
on immediate action utility; it does not run an offline Jev model.

NEAR is a public blockchain supporting smart contracts, programs whose
transactions and state can be checked through the network
\cite{nearprotocol,nightshade2024}.
Boyun's optional testnet registry records an artifact commitment. Testnet is
the network used to test this integration, rather than evidence of a production
deployment. A commitment is a cryptographic fingerprint that binds a later
disclosure to an earlier registration. It is neither encryption nor a proof
that the artifact's contents are correct.

\begin{figure*}[t]
\centering
\begin{tikzpicture}[node distance=7mm and 8mm]
\node[box] (private) {Private artifact\\and fresh secret salt};
\node[box,right=of private] (commit) {Cryptographic commitment\\SHA-256 fingerprint};
\node[box,right=of commit] (chain) {NEAR testnet registry\\Signer and commitment};
\node[box,below=of commit] (open) {Later disclosure\\Recompute and compare};
\draw[arr] (private)--(commit);\draw[arr] (commit)--(chain);
\draw[arr] (private)--(open);\draw[arr] (chain)--(open);
\end{tikzpicture}
\caption{Artifact registration. Original contents and the random salt remain
off chain until deliberately disclosed. The signer and transaction timing
remain public.}
\label{fig:integrity}
\end{figure*}

Algorithm~\ref{alg:commit} describes registration using SHA-256, a hash function
that produces a fixed-size fingerprint. The secret salt is fresh random data;
it prevents a public commitment from being directly compared with hashes of
guessed scenario contents. A fixed purpose label separates this use of hashing
from other protocols. Explicit lengths make field boundaries unambiguous.
Without the salt, predictable private inputs could be tested by guessing;
without purpose and field boundaries, the commitment format would be less
clearly specified. These measures do not hide the public signer or timing.
A registered hash also does not ensure that the off-chain artifact can later
be retrieved. Nightshade's treatment of data availability distinguishes
present block-data availability from long-term retention \cite{nightshade2024};
experiment storage and disclosure remain separate responsibilities here.

\begin{algorithm}[t]
\caption{Registering and opening an artifact commitment}
\label{alg:commit}
\footnotesize
\begin{algorithmic}[1]
\Require Exact artifact bytes; fixed commitment-purpose label
\State Hash the artifact with SHA-256
\State Draw a fresh secret salt of 32 random bytes
\State Encode the purpose label and a zero-byte separator
\State Append the salt length, salt, hash length, and hash
\Statex \hspace{\algorithmicindent}Encode lengths as 64-bit big-endian integers
\State Hash the encoded sequence with SHA-256
\State Register only the resulting commitment on NEAR
\State Retain the artifact and salt privately
\State To open, recompute from the disclosed artifact and salt
\State Check the finalized transaction and registry entry
\end{algorithmic}
\end{algorithm}

Verification checks the transaction, its final outcome, expected contract code,
and matching registry entry through a selected remote procedure call (RPC)
endpoint. RPC is the interface used to query the network; this check is not
an independently verified light-client proof. Contract upgrade authority and
signer metadata remain part of the trust boundary.

NEAR AI Cloud is an optional hosted model service for world authoring
\cite{nearai}, separate from Jev's participant decisions. Although its
documentation describes verification facilities, the present adapter does
not verify a trusted-execution attestation, a cryptographic statement about
the environment processing a request. DCML describes how an attested connection
to known service code and model-hash checking can support confidential,
verifiable inference \cite{dcml}. A normal compatible API request alone
establishes neither property. Consequently, this integration does not
establish confidential inference for the full workflow.
